Les archives professionnelles concernent la conservation et l’exploitation commerciale des archives de l’ORTF et les archives confiées par la loi à l’INA sur lesquels l’institution a un droit d’exploitation. Ce service s’oppose déjà aux archives du dépôt légal dans leur organisation. Les archives professionnelles fonctionnent en « stock », elles gèrent des stocks, des programmes, des émissions à l’unité, qu’elles exploitent, au contraire du dépôt légal qui a une gestion en « flux », selon les tranches horaires, les grilles de programmes\footcite[][23]{varra2023}. Aujourd’hui, ces deux services ont disparu à la suite d’une réorganisation des services en 2020 qui répartit ces deux services dans la « Direction des Patrimoines »\footnote{(Voir organigrammes de la direction des Patrimoines :~\ref{ann:org},\nameref{ann:org}), p.~\pageref{ann:org}.} mais ils transparaissent toujours de diverses manières au sein du traitement documentaire. Nous avons déjà vu que leurs objectifs et leur public n’était pas les mêmes, mais nous n’avons pas vu de quelle façon cela se répercutait sur le traitement documentaire. Pour illustrer nos propos, prenons pour exemple deux notices d’une même émission\footnote{(Pour consulter les notices voir :~\ref{ann:notices},\nameref{ann:notices}), p.~\pageref{ann:notices}.}, une notice provient de Totem, l'outil d'exploration de la base de données des archives professionnelles, l’autre vient de DLTV, la base de données du dépôt légal. Evidemment, les deux notices présentent des informations de description identiques telles que le titre, la date de diffusion, le genre d’émission, le générique et même le sommaire. Mais les informations sur la diffusion diffèrent. Du côté professionnel, nous trouvons la date, l’heure mais l’information importante du côté dépôt légal est que cette émission est en première diffusion. Cette information a peu de valeur dans le cadre d’une commercialisation mais elle en a une à la fois pour les TGDM chargés de déterminer le type de traitement à appliquer à cette émission et pour les chercheurs de l’INAthèque. Les autres informations de diffusion sont présentes dans la fiche Hyperbase mais elles sont réparties entre différents champs afin de permettre la recherche par ce champ. Les autres données de description le confirment. Dans la fiche des archives professionnelles, nous retrouvons les informations concernant le statut de numérisation, essentiel pour partager ce document, la dévolution, sur laquelle nous allons revenir, et surtout des informations précises sur les fichiers, pour la diffusion de ce document. Toutes ces informations sont absentes de la fiche du dépôt légal. Au contraire, nous trouvons dans la fiche du dépôt légal des informations sur les documents d’accompagnements, pouvant compléter une recherche et les audiences permettant, entre autres, de réaliser des statistiques.
 
En comparant ces deux fiches, il est évident que même actuellement, des années après la réorganisation des services, deux politiques documentaires cohabitent. Pourtant, l’écart entre les deux services patrimoniaux est bien moins important qu’il ne le fut. À la création du dépôt légal, il est nécessaire de trouver de nouveaux locaux pour accueillir ce nouveau service. Le service du dépôt légal investit des locaux à une centaine de mètres des locaux historiques de l’INA à Bry-sur-Marne\footcite[][24]{varra2023}, creusant un fossé entre les deux services, même physiquement. C’est un ressenti évoqué par la documentaliste interrogée, anciennement au dépôt légal, qui souligne l’écart entre les deux services. L’accès aux notices, aux images et aux sons des documents des archives professionnelles étaient compliquées avant leur numérisation car ils n’avaient même pas accès à leur outil de consultation « Totem »\footnote{(Voir entretien d'une documentaliste de l'INA sur la différence entre le dépôt légal et les archives professionnelles~\ref{ann:transcriptions},\nameref{ann:transcriptions}), p.~\pageref{ann:transcriptions}.}. Dans ces circonstances, on peut difficilement imaginer une bonne communication sur les politiques documentaires. Nous pouvons également souligner que différents services de l’INA furent séparés géographiquement. Avant même sa création, la cinémathèque était séparée entre le service « Film d’Actualités » à Cognacq-Jay et « Film Production » aux Buttes Chaumont. Ils seront réunis pour une courte durée à la création de l’INA avant de se séparer à nouveau et d’être réunis à Bry-sur-Marne. Nous pouvons donc imaginer, en conséquence, les nouvelles incohérences générées dans les politiques documentaires. 

Cela dit, le manque de communication n’est pas ignoré, à plusieurs reprises, il sera tenté d'harmoniser les méthodes. Des référentiels, des thésaurus sont mis en place mais nous pouvons plus particulièrement souligner la création de la DDCOLL, direction déléguée des collections, en 2010 qui mis en place une structure transversale « méthodes et qualités » et qui donnera lieu à la réunification des services, évoqués plus haut\footcite[][29]{varra2023}, ainsi que la création du « lac de données » sur lequel nous reviendrons plus tard. 
Ces perturbations d’ordre institutionnel ne trouvent pas leur pareil au sein des institutions patrimoniales ayant pour mission le dépôt légal. Il n’y ni au sein du CNC, ni au sein de la BNF de distinction entre les documents du dépôt légal et les documents commercialisés. C’est une exception que l’on doit au statut d’EPIC (Établissement public à caractère industriel et commercial en France) de l’INA que l'on ne retrouve pas à la BnF ou au CNC. 